\documentclass[11pt]{article}
\usepackage[margin=1in]{geometry}
\usepackage[T1]{fontenc}
\usepackage[utf8]{inputenc}
\usepackage{lmodern}
\usepackage{array}
\usepackage{tabularx}
\usepackage{booktabs}
\usepackage{enumitem}
\usepackage{xurl}
\usepackage[hidelinks]{hyperref}
\usepackage{titlesec}

\setlength{\parindent}{0pt}
\setlength{\parskip}{6pt}
\setlist[itemize]{leftmargin=1.5em}
\setlist[enumerate]{leftmargin=1.5em}
\titleformat{\section}{\large\bfseries}{}{0pt}{}
\titleformat{\subsection}{\normalsize\bfseries}{}{0pt}{}
\pagestyle{plain}

\title{Agile Requirements and BABOK\\\large A Practical Alignment Handout}
\author{}
\date{\today}

\begin{document}
\maketitle

\section{Purpose}

This handout explains how continuous requirements work in an Agile software engineering lifecycle and how the \emph{A Guide to the Business Analysis Body of Knowledge} (BABOK Guide) supports that way of working. The core point is straightforward: Agile treats requirements as continuously refined work, while BABOK provides the analysis disciplines, tasks, and techniques needed to manage that refinement in a controlled and repeatable way.\cite{scrumguide,babokmain,babokstructure}

\section{Why Agile Requirements Are Continuous}

In Scrum, requirements are not gathered once and then frozen. The Product Backlog is defined as an emergent, ordered list of what is needed to improve the product. Backlog refinement is described as an ongoing activity that breaks down and further defines backlog items by adding detail, order, and size. Scrum also states that items selected for a Sprint are discussed during Sprint Planning and may continue to be clarified as the team learns.\cite{scrumguide}

In practice, this means requirements work happens at several recurring points:

\begin{itemize}
  \item \textbf{Initial conception and planning:} establish product goals, high-level scope, constraints, and key stakeholder needs.
  \item \textbf{Backlog refinement:} elaborate user stories, acceptance criteria, dependencies, and priorities.
  \item \textbf{Sprint planning:} select the most ready items for near-term delivery and confirm the work needed to complete them.
  \item \textbf{During development and review:} clarify open questions, validate assumptions, and adapt based on demonstrations, feedback, and newly discovered information.
\end{itemize}

This continuous cadence reflects three Agile realities: requirements are progressively elaborated, feedback arrives quickly, and change is expected rather than treated as an exception.\cite{scrumguide}

\section{What BABOK Adds}

The BABOK Guide is IIBA's foundational standard for business analysis. IIBA describes it as a guide to business analysis knowledge areas, tasks, underlying competencies, techniques, and perspectives. The guide is organized around six knowledge areas and five perspectives, including an Agile perspective.\cite{babokmain,babokstructure,babokperspectives}

\subsection{Six knowledge areas}

The six BABOK knowledge areas are:\cite{babokstructure,baboktasks}

\begin{itemize}
  \item Business Analysis Planning and Monitoring
  \item Elicitation and Collaboration
  \item Requirements Life Cycle Management
  \item Strategy Analysis
  \item Requirements Analysis and Design Definition
  \item Solution Evaluation
\end{itemize}

\subsection{Relevant BABOK techniques for Agile work}

BABOK also includes techniques that map directly to Agile delivery, including backlog management, prioritization, and user stories. These are especially relevant when requirements are discovered and refined iteratively instead of being fully specified up front.\cite{babokbacklog,babokprioritization,babokuserstories}

\section{How Agile Cadence Maps to BABOK}

\renewcommand{\arraystretch}{1.25}
\begin{tabularx}{\textwidth}{>{\raggedright\arraybackslash}p{0.24\textwidth} >{\raggedright\arraybackslash}p{0.28\textwidth} X}
\toprule
\textbf{Agile point in time} & \textbf{Primary BABOK knowledge areas} & \textbf{What the analysis work looks like} \\
\midrule
Project or product inception & Strategy Analysis; Business Analysis Planning and Monitoring & Define the business need, desired outcomes, change strategy, stakeholder approach, analysis governance, and information management approach. \\
\addlinespace
Backlog refinement & Elicitation and Collaboration; Requirements Analysis and Design Definition; Requirements Life Cycle Management & Elicit detail, shape user stories, define acceptance criteria, model requirements, verify understanding, trace changes, and reprioritize work. \\
\addlinespace
Sprint planning & Requirements Analysis and Design Definition; Requirements Life Cycle Management & Ensure candidate items are sufficiently understood, prioritized, and ready for implementation in the next Sprint. \\
\addlinespace
During development and Sprint Review & Elicitation and Collaboration; Solution Evaluation & Answer emerging questions, confirm intent with stakeholders, inspect results, and recommend adjustments based on evidence from the increment. \\
\bottomrule
\end{tabularx}

\section{A Practical Operating Model}

A practical way to use Agile and BABOK together is to separate \emph{delivery cadence} from \emph{analysis discipline}:

\begin{enumerate}
  \item Use Agile to determine \textbf{when} requirements are examined, refined, built, demonstrated, and changed.
  \item Use BABOK to determine \textbf{how} requirements are planned, elicited, analyzed, prioritized, governed, traced, and evaluated.
\end{enumerate}

That combined model gives teams a way to stay adaptive without becoming ad hoc. Agile preserves responsiveness; BABOK preserves rigor.

\section{What This Means for Teams}

For a delivery team, the alignment can be summarized as follows:

\begin{itemize}
  \item Do not treat requirements gathering as a single gate at project start.
  \item Maintain a living backlog that can absorb learning and change.
  \item Use structured analysis practices to improve requirement quality before work enters implementation.
  \item Revisit priorities and solution fit continuously, not only at the end of the project.
  \item Evaluate outcomes against stakeholder needs after each increment, then feed that learning back into the backlog.
\end{itemize}

\section{Conclusion}

Agile software development assumes that requirements evolve. The BABOK Guide provides the framework for managing that evolution through defined knowledge areas, tasks, and techniques. Used together, they support a delivery model in which requirements are continuously refined while still being governed with discipline and clear analytical intent.\cite{scrumguide,babokmain,baboktasks}

\section*{References}
\begin{thebibliography}{9}

\bibitem{scrumguide}
Ken Schwaber and Jeff Sutherland,
\emph{The Scrum Guide}, Scrum Guides,
\url{https://scrumguides.org/scrum-guide.html}.

\bibitem{babokmain}
International Institute of Business Analysis,
\emph{A Guide to the Business Analysis Body of Knowledge}, IIBA,
\url{https://www.iiba.org/knowledgehub/business-analysis-body-of-knowledge-babok-guide/}.

\bibitem{babokstructure}
International Institute of Business Analysis,
\emph{Structure of the BABOK Guide}, IIBA,
\url{https://www.iiba.org/knowledgehub/business-analysis-body-of-knowledge-babok-guide/1-Introduction/1-4-structure-of-the-babok-guide/}.

\bibitem{baboktasks}
International Institute of Business Analysis,
\emph{Tasks}, IIBA,
\url{https://www.iiba.org/knowledgehub/business-analysis-body-of-knowledge-babok-guide/tasks/}.

\bibitem{babokperspectives}
International Institute of Business Analysis,
\emph{BABOK Applied}, IIBA,
\url{https://www.iiba.org/knowledgehub/babok-applied/}.

\bibitem{babokbacklog}
International Institute of Business Analysis,
\emph{Backlog Management}, IIBA,
\url{https://www.iiba.org/knowledgehub/business-analysis-body-of-knowledge-babok-guide/10-techniques/10-2-backlog-management/}.

\bibitem{babokprioritization}
International Institute of Business Analysis,
\emph{Prioritization}, IIBA,
\url{https://www.iiba.org/knowledgehub/business-analysis-body-of-knowledge-babok-guide/10-techniques/10-33-prioritization/}.

\bibitem{babokuserstories}
International Institute of Business Analysis,
\emph{User Stories}, IIBA,
\url{https://www.iiba.org/knowledgehub/business-analysis-body-of-knowledge-babok-guide/10-techniques/10-48-user-stories/}.

\end{thebibliography}

\end{document}



